\section{Integrable unipotent majorants}
\label{sec:xi}

The Fourier calculation controls oscillation. Absolute convergence of the real Bessel integrals also requires decay along the unipotent radical, including directions of the flag that the character does not suppress.

Write \(\operatorname{Del}_l(g)\) for the square root of the Gram determinant of the first \(l+1\) vectors in the translated null frame, with the Euclidean inner product. For small positive defects \(e_l\) from the reference exponents \(s_l^0\), the positive function on the real unitary group used in the orbital integral is 
\[
 \Xi(g)=\prod_l\operatorname{Del}_l(g)^{-(s_l^0-e_l)}.
\]
 Proposition~\ref{prop:xi-group-majorant} uses this function for both translations and establishes integrability on the unipotent radical. The exponent for the
group norm in this calculation can be taken to be
\(A\ge 2(k+r)(r+1)+1\), where the unitary group has signature
\((k+r,r+1)\). 

The following argument proves existence at each rank. The constants in this existence argument may
depend on the rank and the permitted exponent defects.
This suffices for absolute convergence of the local integrals
and the exchanges of order of integration used in the Bessel
normalization and local amplitude estimates of
Sections~\ref{sec:bessel-periods} and~\ref{sec:local-amplitudes}.

\subsection{Heisenberg integrals}

Successive integration along a null flag removes one Heisenberg group at a time. Its central coordinate and its two isotropic coordinates enter the Gram minors differently.

Use variables \(a,b\in\C^p\), \(c\in\C^m\) and \(\beta\in\R\).
With indices \(0\le l\le p\) set
\[
 \begin{split}
 L_l&=1+\sum_{j<l}|a_j|^2,\qquad
 R_l=L_p+\|c\|^2+\sum_{j\ge l}|b_j|^2,\\
 B_l&=-\|c\|^2/2+i\beta-\sum_{j\ge l}b_j\overline{a_j},
 \qquad D_l^2=L_lR_l+|B_l|^2.
 \end{split}
\]
The critical list is \(s_l^0=2\) for \(l<p\) and
\(s_p^0=m+1\).

\begin{lemma}[Heisenberg integrability]\label{lem:heisenberg-integrability}
The integral with exponents \(-s_l^0-\nu_l\) is finite
if all prefix sums of \(\nu_l\) are positive and every
proper prefix plus twice its tail is positive.
For \(p\ge1\), the integral of
\[
 |b_0|^{-\gamma}\prod_l D_l^{-s_l^0+e_l}
\]
is finite if \(0<\gamma\le1\), \(e_l>0\), and
\(e_0+2\sum_{l>0}e_l<\gamma/2\).
\end{lemma}

\begin{proof}
\proofstep{The unsigned integral.}
The integral of
\(\prod_l D_l^{-s_l^0-\nu_l}\) is finite when the
prefix sums of \(\nu\) are positive and every proper
prefix plus twice its tail is positive. If \(p=0\),
\(D_0^2=(1+\|c\|^2/2)^2+\beta^2\), so integrating \(\beta\)
leaves a constant times \((1+\|c\|^2/2)^{-m-\nu_0}\).

For \(p\ge1\) put \(R=1+|a_0|^2\). The first square completes as
\[
 D_0^2=R|b_0-a_0B_1/R|^2+D_1^2/R .
\]
The first prefix condition gives \(\nu_0>0\). Translation
of \(b_0\), followed by polar coordinates, gives the exact identity
\[
 \int_{\mathbb C}D_0^{-2-\nu_0}\,d^2b_0
  =\frac{\pi}{R}\int_{D_1^2/R}^{\infty} t^{-1-\nu_0/2}\,dt
  =\frac{2\pi}{\nu_0}R^{-1+\nu_0/2}D_1^{-\nu_0}.
\]

More generally, with a prefactor
\(R^{-1+\alpha/2}D_1^{-\alpha}\), scale the remaining
complex coordinates by \(\sqrt R\) and \(\beta\) by \(R\).
The remaining \(D_l\) scale by \(R\), while the measure
scales by \(R^{2(p-1)+m+1}\). The reference exponents
\(\sum_{l>0}s_l^0=2(p-1)+m+1\) cancel this Jacobian.
The resulting \(a_0\)-density is
\[
 R^{-1+\alpha/2}\,R^{-\alpha-\sum_{l>0}\nu_l}d^2a_0
   =R^{-1-\alpha/2-\sum_{l>0}\nu_l}d^2a_0 .
\]
The lower-rank integral merges \(\alpha\) into \(\nu_1\).
For \(\alpha=\nu_0\) the integral over \(a_0\in\mathbb C\) is
finite precisely when
\[
             \nu_0+2\sum_{l=1}^p\nu_l>0.
\]
This is the first proper-prefix condition. The reduced list is
\((\nu_0+\nu_1,\nu_2,\ldots,\nu_p)\); its prefix and tail tests
are exactly the remaining tests on the original list. Induction
therefore proves convergence.

\proofstep{A singular weight in the first opposite coordinate.}

The unsigned integral controls the tail at infinity. The character estimate will also insert a singular power of one opposite coordinate, so local integrability at its zero has to be checked.


There is a second integral with all exponents just below
critical. If \(p\ge1\), \(0<\gamma\le1\), \(e_l>0\) and
\(e_0+2\sum_{l>0}e_l<\gamma/2\), then
\[
  \int |b_0|^{-\gamma}\prod_l D_l^{-s_l^0+e_l}<\infty .
\]
Indeed for \(0<e<\gamma<2\) and \(d>0\),
\[
 \int_{\C} (|x-h|^2+d^2)^{-1+e/2}|x|^{-\gamma}\,d^2x
 \le C_{e,\gamma} (|h|^2+d^2)^{-(\gamma-e)/2}.
\]
Rescale to \(|h|^2+d^2=1\).

If \(d\ge1/2\), the
first factor is bounded on bounded sets. Otherwise
the singularities at \(0\) and \(h\) are separated.
The local powers are integrable at either point and the
infinite tail has exponent \(-2+e-\gamma\).
The square completion now gives the bound
\(C R^{-1+\gamma/2}D_0(b_0=0)^{-(\gamma-e_0)}\).
Put \(\alpha=\gamma/2-e_0\). Since \(D_0(b_0=0)\ge
D_1/\sqrt R\) and \(D_0(b_0=0)\ge|\beta-\beta_*|\),
this bound is at most
\[
 C R^{-1+\gamma/2+\alpha/2}
                  D_1^{-\alpha}|\beta-\beta_*|^{-\gamma/2}.
\]
The remaining density in \(\beta\) is a product of
negative powers of \((A_l^2+(\beta-\beta_l)^2)^{1/2}\)
with all widths \(A_l\ge R\).

On
\([\beta_*-R,\beta_*+R]\) each square root varies by a
factor at most three. Averaging the density over that
interval and integrating the singular factor bounds the
integral by \(C R^{-\gamma/2}\) times the integral of the
density. The same comparison holds on the complement
since \(|\beta-\beta_*|\ge R\). We are left with the
previous prefactor \(R^{-1+\alpha/2}D_1^{-\alpha}\)
and \(\nu_l=-e_l\) for \(l>0\). Its conditions follow
from \(\alpha>2\sum_{l>0}e_l\).

\end{proof}

\subsection{The character weight}

Inserting the unipotent character introduces a real linear form in addition to the Gram weights. Decay in its translate is needed when the remaining coordinates are integrated.

\begin{proposition}[A character-weighted Heisenberg integral]\label{prop:heisenberg-character}
For the Heisenberg coordinates, minor exponents \(s_l^0-e_l\)
near their reference values, and linear form with nonzero
coefficient \(\eta\) on \(b_0\), the inserted weight
\((1+|2\operatorname{Im}\mathcal L+c_0|)^{-\delta}\)
gives an integral bounded by
\(C|\eta|^{-\gamma}(1+|c_0|)^{-\delta/4}\) in the range 
\[
 p\ge1,\quad 0<\delta\le\tfrac12,\quad 0<\gamma\le\delta/8,\quad e_l>0,\quad e_0+2\sum_{l>0}e_l<\gamma/2.
\]
 The constant
is independent of the coefficients of the linear form and \(c_0\).
\end{proposition}

\begin{proof}
\proofstep{A quadratic sublevel estimate.}

A polynomial of degree at most two on \([1,2]\) with maximum
coefficient modulus \(M>0\) has a sublevel set of measure
at most \(C\sqrt u\) where its modulus is at most \(uM\).
In a sublevel set of positive measure \(v\) choose three
points \(t_i\) separated by at least \(v/4\). Lagrange interpolation gives
\[
 P(t)=\sum_{i=1}^3 P(t_i)\prod_{j\ne i}\frac{t-t_j}{t_i-t_j}.
\]
Each coefficient on the right has magnitude \(O(uM/v^2)\),
so \(v^2\ll u\). Splitting at \(u=M^{-1/2}\) proves
\(\int_1^2(1+|P(t)|)^{-\delta}\,dt\le
C_\delta(1+M)^{-\delta/2}\) for \(0<\delta\le1/2\).

\proofstep{Averaging the linear character.}

Consider \(2\operatorname{Im}\mathcal L+c_0\), with
\(\mathcal L\) linear in the complex coordinates and nonzero
coefficient \(\eta\) on \(b_0\). Rotate all complex
coordinates by \((1+it)^2/(1+t^2)\).
The \(D_l\) and measure are invariant. Write \(\mathcal L=A+iB\) before rotation. After multiplying
by \(1+t^2\), the phase is
\[
     (c_0+2B)+4At+(c_0-2B)t^2 .
\]
Its maximum coefficient is comparable to \(|c_0|+|\mathcal L|\).
The sublevel estimate thus bounds the averaged weight
by \(C(1+|c_0|)^{-\delta/4}(1+|\mathcal L|)^{-\delta/4}\).
Average next over \(a_0\mapsto ta_0,b_0\mapsto b_0/t\)
for \(t\in[1,2]\). The measure is invariant, and each \(D_l\)
is comparable to its original value by fixed factors.

Write the resulting linear form as \(tA'+t^{-1}\eta b_0+C_0\).
Multiplication by \(t\) gives the quadratic
\[
          t^2 A'+t C_0+\eta b_0 .
\]
Its maximum coefficient is at least \(|\eta b_0|\), so the
averaged weight is at most \(C(1+|\eta b_0|)^{-\delta/8}\).
For \(0<\gamma\le\delta/8\) apply the singular integral
above. This gives
\[
 \int \prod_l D_l^{-s_l^0+e_l}
       (1+|2\operatorname{Im}\mathcal L+c_0|)^{-\delta}
 \ \le\ C|\eta|^{-\gamma}(1+|c_0|)^{-\delta/4}.
\]
The constant is uniform in the coefficients of the
linear form and in \(c_0\).

\end{proof}

\begin{proposition}[Character-weighted convergence on the flag unipotent]\label{prop:xi-flag-unipotent}
In the flag coordinates and with the Bessel character under
consideration, the character-weighted unipotent integral with
minor exponents near \(s_l^0\) is finite for sufficiently small
positive defects. At fixed rank and for \(0<\delta\le1/2\), the induction chooses positive defects at each radical. Its character exponent after one radical and its permitted loss there are 
\[
 \delta_{\rm next}=\delta/4,\qquad 0<\gamma\le\delta/8,
\]
 by Proposition~\ref{prop:heisenberg-character}. The measure is product Lebesgue measure in the successive Heisenberg coordinates.
\end{proposition}

\begin{proof}
\proofstep{Identifying the unipotent integrals.}

For the group interpretation remove the first hyperbolic
pair \(e,f\), with form \(J\) on the remaining space.
The Heisenberg radical has matrices
\[
 h(x,t)=
 \begin{pmatrix}
 1&-x^*J&-x^*Jx/2+it\\0&1&x\\0&0&1
 \end{pmatrix}.
\]
Multiplying the two matrices gives
\[
 h(x,t)h(y,u)=h\bigl(x+y,t+u-\operatorname{Im}(x^*Jy)\bigr).
\]
Left or right translation has triangular Jacobian with
diagonal entries one. Thus \(dx\,dt\) is Haar in these
coordinates. The smaller unipotent acts linearly on \(x\)
by a determinant-one unitary matrix and fixes the central
parameter \(t\), so the successive product measure is Haar.
In the opposite null frame \(v_j\), its Euclidean dual
and the remaining positive directions, the variables are
\[
   x=\sum_j(a_jJv_j+b_jv_j)+c,\qquad
   \beta=t+\operatorname{Im}\sum_j b_j\bar a_j .
\]
The shear is translation in the central coordinate, so it
preserves measure. The Gram determinant of the first
column \(h(x,t)f\) and the next \(l\) null columns
\(h(x,t)v_j=v_j-\bar a_j e\) is \(D_l^2\).

In fact the latter columns have Gram determinant
\(L_l\). Subtracting their \(b_j\)-multiples from
the first column leaves middle and last squared
length \(R_l\) and top coefficient \(B_l\).
The Gram matrix of the last \(l\) columns is an identity
plus the rank-one matrix from \((a_j)_{j<l}\). If \(e\) is
the top unit vector, its squared projection on their orthogonal
complement is \(1-\|(a_j)_{j<l}\|^2/L_l=1/L_l\).
The Schur complement for the first column thus gives
\[
       D_l^2=L_l\bigl(R_l+|B_l|^2/L_l\bigr).
\]

\proofstep{Induction in the flag rank.}

The one-step estimate retains dependence on the translate of the character. That retained decay permits integration over the next layer of the flag.


For a fixed smaller element \(m_0\), Euclidean
Gram--Schmidt writes its translated opposite frame
as a compact translate times a triangular matrix
\(P\). For the positive and negative projections of an orthonormal
isotropic frame, the two Gram matrices sum to the identity and
have zero difference. Thus
\[
    G_++G_-=I,\qquad G_+-G_-=0,\qquad G_+=G_-=I/2 .
\]
The compact group acts transitively on these pairs of scaled
orthonormal frames. The triangular multiplier at each stage
is the corresponding principal Euclidean Gram
square root \(\operatorname{Del}_l(m_0)\).
In the new Heisenberg variables the character has
constant term equal to the smaller character phase
and linear part pairing \(x\) with \(m_0 w_*\) after
the compact frame change, for the first smaller
Bessel vector \(w_*\). Thus
\[
 |P_{11}|\,|\eta|=|\langle v_1,w_*\rangle_J|,
 \qquad |P_{11}|=\operatorname{Del}_0(m_0).
\]
The pairing is the dual null pairing if the smaller
rank is positive.

In rank zero \(w_*\) is negative
and a null perpendicular space is nonnegative.
The inserted coefficient is therefore nonzero.
By Proposition~\ref{prop:heisenberg-character}, the integration
over the first radical inserts \(|\eta|^{-\gamma}\) and leaves
a smaller character weight of exponent \(\delta/4\).
The displayed identity for \(\eta\) makes this the insertion
of \(\operatorname{Del}_0(m_0)^\gamma\), times a constant.
Thus only the first smaller minor exponent is reduced, by
\(\gamma\).
If the smaller rank permits defects less than \(e'\),
choose \(0<\gamma\le\min(\delta/8,e'/2)\) and take
the new defects less than \(\min(e'/2,\gamma/(4p+4))\).
The induction starts with the trivial unipotent.

\end{proof}

\begin{construction}
\begin{tabular}{@{}ll@{}}\toprule
Columns at one Heisenberg step & Euclidean Gram data\\\midrule
\(h(x,t)v_j=v_j-\bar a_j e\) & determinant \(L_l\)\\
\(h(x,t)f\), projected off these columns &
squared length \(R_l+|B_l|^2/L_l\)\\
All \(l+1\) columns & determinant \(D_l^2=L_l R_l+|B_l|^2\)\\\bottomrule
\end{tabular}
\caption{Computing one flag minor by orthogonal projection. The rank-one correction to the Gram matrix of the last \(l\) columns contributes the denominator \(L_l\).}\label{fig:heisenberg-gram-minor}
\end{construction}

\subsection{The group bound}

The one-step integrations control the successive unipotent coordinates. A uniform estimate on the real group must also allow a translate of the null frame on either side.

Let \(a(g)=\|g\|_{\mathrm{Frob}}\) and write \(F_s(g)\)
for the product of \(-s_l\)-th powers of the absolute
values of the leading Hermitian minors of \(g\) on the
opposite null flag. Take \(s_l=s_l^0-e_l\) with small
positive defects. For a bi-compact-invariant weight,
translating \(F_s\) on either side translates the
associated null frame. Gram--Schmidt and compact
invariance then factor out the products of the
principal Euclidean Gram square roots to the powers
\(-s_l\). This gives \(\Xi(x)\Xi(y)\) for translation
by \(x,y\), with \(\Xi=\prod_l\operatorname{Del}_l^{-s_l}\).

\begin{proposition}[Integrability for the two-sided majorant]\label{prop:xi-group-majorant}
With the small positive defects of the Heisenberg argument,
the group of signature \((k+r,r+1)\) satisfies 
\[
 \int_G F_s(g)a(g)^{-A}\,dg<\infty\qquad\bigl(A\ge2(k+r)(r+1)+1\bigr).
\]
 The resulting
two-sided weight is bounded by a fixed constant times
\(\Xi(x)\Xi(y)\), with the corresponding absolute unipotent
integrals in the same normalization.
\end{proposition}

\begin{proof}
\proofstep{The full-measure cell.}

We check \(\int_G F_s(g)a(g)^{-A}dg<\infty\).
With \(n\) the dimension after removing \(e,f\), write
the full-measure cell as
\(h(x,t)^{-1}w(c)\operatorname{lev}(g_0)h(y,u)\).
Here \(w(c)e=cf,\ w(c)f=\bar c^{-1}e\).
In the first-column coordinates, scaling the \(n\) unrestricted
complex entries and the complex \(e\)-coefficient contributes
\(|c|^{2(n+1)}\). The real null equation contributes
\(|c|^{-2}\), leaving density
\[
   |c|^{2n}\,d^2c\,dx\,dt .
\] The stabilizer of
the first column has product Haar measure on the lower
group and Heisenberg radical. Each leading Hermitian
minor factors as the product of \(c\) and its smaller
minor, by first row and column elimination. The Haar
integral thus has a factor \(|c|^{2n-\sum_l s_l}\).

\proofstep{Two bounds on the right radical.}

Put \(d=2n+1\). Integration over the right radical at
fixed \(c,g_0\) satisfies
\[
 E(c,g_0)=\int a(w(c)\operatorname{lev}(g_0)h(y,u))^{-A}dy\,du
 \le C_A\min(|c|^{-d},|c|^{-1})a(g_0)^{d-A}.
\]
The squared norm bounds the two expressions
\[
   a(g_0)^2+|c|^2(\|y\|^2+u^2),\qquad
   a(g_0)^2+\|g_0y\|^2+|c|^2u^2.
\]
For the first use \((y,u)\mapsto |c|(y,u)/a(g_0)\).
For the second use \((y,u)\mapsto(g_0y,|c|u)/a(g_0)\), where
\(|\det_{\R}g_0|=1\). The inverse Jacobians give
\(|c|^{-d}a(g_0)^d\) and \(|c|^{-1}a(g_0)^d\), respectively.

\proofstep{Induction after fixing the first column.}

Let \(D_0=\|h(x,t)^{-1}f\|\). Use a compact element
carrying \(f\) to \(D_0^{-1}h(x,t)^{-1}f\).
If the chosen compact element is \(k_0\), its removal gives
the first-column computation
\[
    w(cD_0)^{-1}k_0^{-1}\bigl(c h(x,t)^{-1}f\bigr)
                    =w(cD_0)^{-1}(cD_0f)=e .
\]
The remaining factor is a lower group element times
a right radical. Its smaller Euclidean Gram square
roots are \(D_l/D_0\) for \(l>0\), by taking the
Gram determinant after projection off the first
column.

Translate in the right radical and apply
the compact-transfer identity on the lower group.
This bounds the lower integral of the norm by
the lower induction at exponent \(A-d\), multiplied
by \(\prod_{l>0}(D_l/D_0)^{-s_l}\) and the right
radical estimate at \(cD_0\).
The threshold \(2(k+r)(r+1)+1\) drops by \(d+1\)
when \(r\) decreases by one, as required.
For \(\sigma=\sum_l s_l\) we have \(1<\sigma<2n+1\).
In the bound for the right radical let \(c'=cD_0\). The
weighted integral in \(c\) is
\[
 D_0^{-(2n-\sigma+2)}
   \int_{\C}|c'|^{2n-\sigma}\min(|c'|^{-d},|c'|^{-1})\,d^2c' .
\]

Since \(d=2n+1\), the radial part of the last integral is
a constant times
\[
    \int_0^1 t^{2n-\sigma}\,dt
        +\int_1^\infty t^{-\sigma}\,dt .
\]
It is finite exactly for \(1<\sigma<2n+1\).
After multiplying by \(\prod_{l>0}(D_l/D_0)^{-s_l}\)
the remaining Heisenberg integral has defects
\(\nu_0=e_0+2\sum_{l>0}e_l\) and \(\nu_l=-e_l\) for \(l>0\).
Thus each prefix has positive defect sum,
\[
 \sum_{l=0}^j\nu_l=e_0+2\sum_{l>0}e_l-\sum_{0<l\le j}e_l>0.
\]

All those sums are small when the original defects are small.

\proofstep{The last group integral.}

At this point both radical integrations have bounds in the remaining variables. The group integral is finite only if the exponents left after both integrations still give decay.


The group base is \(r=0\) with smaller compact
\(U(k-1)\). For \(k\ge2\) the preceding \(c\)-integral
leaves the single tail with exponent \(k+e_0\).
For \(k=1\) the two extreme entries give the bound
\[
 E(c,1)\le C_A |c|^{-1}(|c|^2+|c|^{-2})^{(1-A)/2}.
\]
The \(c\)-integral then converges for \(A\ge3\) and
\(0<s_0<1\). The remaining central variable has
exponent \(2-s_0>1\).

Finally each leading minor is bounded by a constant
times \(a(g)^{l+1}\), so \(a(g)^{-A}\le C F_s(g)\)
at these exponents. The shifted group integral
therefore supplies the required product majorant.
The integrable unipotent tails follow from the
character bound above. Compact factors in the two
definite stabilizers normalize the unipotent and
preserve its measure and character. Their smallest
singular values have positive lower bounds, giving
fixed lower bounds on the translated Gram square
roots. The same \(\Xi\) works for the group and
unipotent integrals.
\end{proof}
